\section*{الفصل \ref{ch:g10:soil-formation} --- التربة، مورد بطيء التكوّن}

\begin{solution}{exo:g10:soil-formation:1}
\omterm{def:g10:soil-formation:chemical-weathering}{التجوية الكيميائية} هي تحوّل معادن الصخر بتفاعلات مع الماء والغازات التي يحملها
(ثاني أكسيد الكربون، الأكسجين): فبعض المعادن يذوب، وبعضها يصير معادن جديدة مستقرة
عند السطح. \omterm{def:g10:soil-formation:hydrolysis}{والحلمهة} هي \omterm{def:g10:soil-formation:chemical-weathering}{التجوية الكيميائية} لسيليكات بالماء وأيونات الهيدروجين فيه:
تذهب أيونات الفلز في المحلول ويُعاد بناء بقية المعدن معدنًا
طينيًا.
\end{solution}

\begin{solution}{exo:g10:soil-formation:2}
\omterm{def:g10:soil-formation:weathering-profile}{السابروليت} صخر متحلل في مكانه: لين وترابي، يحتفظ بحبيبات صخره الأم وأشرطته وبنيته،
ويتصل نحو الأسفل، دون انقطاع، بالصخر السليم. أما التربة المتكونة في رمل نهري فتستقر
على مادة جُلبت من مكان آخر: حبيباتها مستديرة ومفروزة، ولا تُظهر شيئًا من بنية الصخر
الذي تحتها.
\end{solution}

\begin{solution}{exo:g10:soil-formation:3}
تصير الفلسبارات طينًا (كاولينيت)، محررةً البوتاسيوم أو الصوديوم أو الكالسيوم
والسيليكا في الماء؛ وتصير الميكا السوداء طينًا وأكاسيد حديد؛ ويبقى الكوارتز دون
تغير، رملًا.
\end{solution}

\begin{solution}{exo:g10:soil-formation:4}
يتأكسد حديد المعادن الداكنة، \ce{Fe^2+}، بالأكسجين الذي يحمله الماء فيصير أكاسيد
حديد(III) غير ذائبة: الغوثيت الأصفر البني والهيماتيت الأحمر، اللذين يغلّفان الحبيبات
ويلوّنان المادة كلها.
\end{solution}

\begin{solution}{exo:g10:soil-formation:5}
\qty{0.017}{mm} في السنة و\qty{1.537}{mm} في السنة: $1.537/0.017 \approx 90$
مرة أكبر.
\end{solution}

\begin{solution}{exo:g10:soil-formation:6}
الكالسيوم: 1 و1. الكربون: $1 + 1 = 2$ و2. الأكسجين: $3 + 2 + 1 = 6$
و$2 \times 3 = 6$. الهيدروجين: 2 و2. الشحنة: 0، و$+2 - 2 = 0$.
\end{solution}

\begin{solution}{exo:g10:soil-formation:7}
لا ينتج ذوبان الكالسيت إلا أيونات ذائبة، \ce{Ca^2+}
و\ce{HCO3-}: فكل شيء يُحمل بعيدًا ولا يترك الحجر الجيري النقي أي بقايا صلبة. أما
\omterm{def:g10:soil-formation:hydrolysis}{حلمهة} الفلسبار فتنتج مادة صلبة، الكاولينيت، تبقى في مكانها؛ وكوارتز الغرانيت لا يكاد
يتفاعل. ولذلك يترك الغرانيت طينًا ورملًا.
\end{solution}

\begin{solution}{exo:g10:soil-formation:8}
تحمل الرقائق الطينية والدبال شحنات سالبة تمسك أيونات البوتاسيوم الموجبة \ce{K+} على
أسطحها، فلا يحملها ماء المطر المار. ويحرر الجذر أيونات هيدروجين تحل محل البوتاسيوم
على الرقائق؛ فيأخذ الجذر البوتاسيوم المتحرر في ماء
التربة.
\end{solution}

\begin{solution}{exo:g10:soil-formation:9}
تحتاج \omterm{def:g10:soil-formation:chemical-weathering}{التجوية الكيميائية} إلى الماء، الوفير في المناخات الرطبة والشحيح في الجافة
(والمتجمد في الباردة)؛ وتفاعلاتها تجري أسرع في الدفء. ولا تمضي التفاعلات بعيدًا بما
يكفي لإزالة كل ما يذوب تقريبًا وترك أكاسيد الحديد والألومنيوم إلا مع الحرارة والماء
الوفير، على مدى زمن طويل.
\end{solution}

\begin{solution}{exo:g10:soil-formation:10}
$10/0.017 \approx 590$ سنة لتكوين \qty{1}{cm}؛ $300/0.017 \approx 17\,600$
سنة لتكوين \qty{30}{cm}؛ $1000/0.017 \approx 59\,000$ سنة لتكوين \qty{1}{m}.
\end{solution}

\begin{solution}{exo:g10:soil-formation:11}
الحجم: $10\,000 \times 0.0015 = \qty{15}{m^3}$؛ والكتلة: $15 \times 1200 =
\qty{18000}{kg}$، أي \num{18}~طنًا لكل هكتار في السنة.
\end{solution}

\begin{solution}{exo:g10:soil-formation:12}
$46\,000/0.017 \approx 2.7$ مليون سنة. فمثل هذا \omterm{def:g10:soil-formation:weathering-profile}{السابروليت} السميك يحتاج إلى مشهد
طبيعي بقي مستقرًا، دون أن يُكشط غطاؤه المتجوّي، ملايين السنين.
\end{solution}

\begin{solution}{exo:g10:soil-formation:13}
الغوثيت: \ce{4 Fe^2+ + O2 + 6 H2O -> 4 FeO(OH) + 8 H+}: الحديد 4 و4،
والأكسجين $2 + 6 = 8$ و$4 \times 2 = 8$، والهيدروجين 12 و$4 + 8 = 12$، والشحنة
$+8$ و$+8$. والهيماتيت: \ce{4 Fe^2+ + O2 + 4 H2O -> 2 Fe2O3 + 8 H+}: الحديد 4
و4، والأكسجين $2 + 4 = 6$ و6، والهيدروجين 8 و8، والشحنة $+8$ و$+8$.
\end{solution}

\begin{solution}{exo:g10:soil-formation:14}
الفقدان الصافي $0.082 - 0.017 = \qty{0.065}{mm}$ في السنة؛ $300/0.065 \approx 4600$
سنة. وتحت الحراثة التقليدية، $300/1.52 \approx 200$ سنة: أي أقصر بأكثر من
عشرين مرة.
\end{solution}

\begin{solution}{exo:g10:soil-formation:15}
يحتوي الطن على \qty{300}{kg} من الأورثوكلاز. وكل كيلوغرام يعطي
\qty{464}{g} من الكاولينيت: $300 \times 0.464 \approx \qty{139}{kg}$. والبوتاسيوم:
$300 \times 0.1405 \approx \qty{42}{kg}$.
\end{solution}

\begin{solution}{pb:g10:soil-formation:1}
\textbf{1.} $4 \times 0.3048 \approx \qty{1.22}{m}$.

\textbf{2.} $0.5 \times 30.48 \approx \qty{15}{cm}$.

\textbf{3.} الدبال، المادة العضوية الداكنة التي تصنعها المحللات من النباتات
الميتة.

\textbf{4.} أكاسيد الحديد، الغوثيت والهيماتيت، المتكونة بأكسدة حديد(II) معادن
المادة الأم.

\textbf{5.} لا. ففي الأسفل يتصل المقطع بصخر متجوٍّ، هو \omterm{def:g10:soil-formation:weathering-profile}{السابروليت} (أو مادة متجوّية
أخرى)، ثم \omterm{def:g10:soil-formation:weathering-profile}{بالصخر الأم}
السليم.

\textbf{6.} \ce{2 KAlSi3O8 + 2 H+ + 9 H2O -> Al2Si2O5(OH)4 + 2 K+ + 4 H4SiO4}.

\textbf{7.} $1000/278.33 \approx \qty{3.59}{mol}$.

\textbf{8.} $3.59/2 \approx \qty{1.80}{mol}$ من الكاولينيت، أي $1.80 \times 258.16
\approx \qty{464}{g}$.

\textbf{9.} أيونات البوتاسيوم والسيليكا الذائبة: يحملها الماء إلى الأنهار، وتحملها
الأنهار إلى البحر.

\textbf{10.} لا يكاد يتفاعل: فتبقى حبيبات الكوارتز رملًا.

\textbf{11.} $1219/0.017 \approx 72\,000$ سنة.

\textbf{12.} $1.537 - 0.017 = \qty{1.52}{mm}$ في السنة.

\textbf{13.} $1219/1.52 \approx 800$ سنة.

\textbf{14.} $72\,000/800 = 90$: فالمقطع يُفقد أسرع بتسعين مرة مما
تكوّن.

\textbf{15.} لا: $0.013 < 0.017$، فالتعرية أبطأ قليلًا من الإنتاج والتربة تحافظ على
نفسها.

\textbf{16.} $T = S/(E - P)$.

\textbf{17.} $300/1.52 \approx 200$ سنة.

\textbf{18.} $300/(0.082 - 0.017) = 300/0.065 \approx 4600$ سنة.

\textbf{19.} التربة التي تُفقد في بضعة قرون تحتاج إلى عشرات آلاف السنين لتتكون من
جديد: فعلى مقياس أعمار البشر والمجتمعات، ما يُفقد لا يُعوَّض.

\textbf{20.} $1000/1.52 \approx 660$ سنة: \textbf{يُفقد متر من التربة
في نحو 660 سنة، أي بضعة قرون، تحت الحراثة التقليدية}.
\end{solution}
